\section*{Bab \ref{ch:b3:point-groups} --- Simetri dan Grup Titik}

\begin{solution}{exo:b3:point-groups:1}
Etena: tiga $C_2$ yang saling tegak lurus, tiga bidang, $i$: $D_{2h}$. \ce{XeF4}
(segi empat datar): $C_4$, empat $C_2 \perp C_4$, $\sigma_h$, $2\sigma_v$, $2\sigma_d$,
$i$, $S_4$: $D_{4h}$. \ce{PCl5} (bipiramida trigonal): $C_3$, tiga $C_2$, $\sigma_h$,
$3\sigma_v$, $S_3$: $D_{3h}$. 1,3,5-Triklorobenzena: unsur-unsur yang sama dengan
\ce{BF3}: $D_{3h}$.
\end{solution}

\begin{solution}{exo:b3:point-groups:2}
Etana bersilang: $C_3$, tiga $C_2$ yang tegak lurus, tiga $\sigma_d$, $i$, $S_6$:
$D_{3d}$. Eklips: $C_3$, tiga $C_2$, $\sigma_h$, tiga $\sigma_v$: $D_{3h}$.
Sikloheksana kursi: $D_{3d}$.
\end{solution}

\begin{solution}{exo:b3:point-groups:3}
\ce{SO3} ($D_{3h}$): tidak. \ce{SO2} ($C_{2v}$): ya. \ce{PCl3} ($C_{3v}$): ya.
\ce{XeF4} ($D_{4h}$): tidak. \ce{CH2Cl2} ($C_{2v}$): ya. cis-1,2-dikloroetena
($C_{2v}$): ya; trans ($C_{2h}$): tidak.
\end{solution}

\begin{solution}{exo:b3:point-groups:4}
$C_2(z) = \mathrm{diag}(-1,-1,1)$, $\chi = -1$; $i = \mathrm{diag}(-1,-1,-1)$,
$\chi = -3$; $\sigma_h(xy) = \mathrm{diag}(1,1,-1)$, $\chi = 1$.
\end{solution}

\begin{solution}{exo:b3:point-groups:5}
Dengan matriks-matriks diagonal dari latihan sebelumnya: $C_2i = \sigma_h$,
$C_2\sigma_h =
i$, $i\sigma_h = C_2$, setiap operasi merupakan inversnya sendiri,
dan setiap hasil kali bersifat komutatif. Grup itu komutatif, sehingga
$X^{-1}AX = A$ untuk semua $X$: setiap operasi merupakan satu kelas tersendiri
(empat kelas, empat \omterm{def:b3:point-groups:irrep}{representasi tak tereduksi} berdimensi 1).
\end{solution}

\begin{solution}{exo:b3:point-groups:6}
Ambil $\sigma_v(1)$ sebagai bidang $xz$. Titik $(1,0)$ dibawa oleh $C_3$ ke $(-\frac12,
\frac{\sqrt3}2)$, lalu oleh $\sigma_v(1)$ ke $(-\frac12,-\frac{\sqrt3}2)$; dengan urutan
sebaliknya titik itu pindah ke $(1,0)$ lalu ke $(-\frac12,\frac{\sqrt3}2)$. Kedua
hasil kali itu adalah pencerminan terhadap dua bidang yang berbeda
($\sigma_v(3)$ dan $\sigma_v(2)$): grupnya tidak komutatif, itulah sebabnya
pencerminan-pencerminannya membentuk satu kelas yang beranggota tiga.
\end{solution}

\begin{solution}{exo:b3:point-groups:7}
Bifenil yang terpuntir mempertahankan $C_2$ sepanjang ikatan antarcincin dan
dua $C_2$ yang tegak lurus padanya, yang membagi dua sudut antara bidang-bidang
cincin; semua bidang dan \omterm{def:b3:point-groups:inversion}{pusat inversi} bentuk datar atau bentuk tegak lurus
hilang. Dengan tiga $C_2$ yang saling tegak lurus dan tanpa unsur tak sejati,
grupnya adalah $D_2$: molekul itu kiral (kedua bentuk terpuntirnya adalah
enantiomer, yang dapat dipisahkan jika substituen orto yang besar mencegah
rotasi).
\end{solution}

\begin{solution}{exo:b3:point-groups:8}
$d_{z^2}$ dan $d_{x^2-y^2}$: $a_1$ ($x^2$, $y^2$, $z^2$ adalah $A_1$); $d_{xy}$:
$a_2$; $d_{xz}$: $b_1$; $d_{yz}$: $b_2$.
\end{solution}

\begin{solution}{exo:b3:point-groups:9}
$1 + 1 + 4 + 9 + 9 = 24 = h$. $\sum g\,\chi_{A_1}\chi_{T_2} = 1\cdot3 + 8\cdot0 + 3(-1)
+ 6(-1) + 6(1) = 0$.
\end{solution}

\begin{solution}{exo:b3:point-groups:10}
Pada bidang yang tegak lurus garis potong, pencerminan terhadap sebuah garis
dengan sudut $\alpha$ memetakan sudut kutub $\phi$ ke $2\alpha - \phi$. Dua
pencerminan, pada $\alpha$ lalu $\alpha + \theta$: $\phi \to 2\alpha - \phi \to 2(\alpha + \theta) - (2\alpha - \phi) =
\phi + 2\theta$, yaitu rotasi sebesar $2\theta$ terhadap garis itu. Jadi dua
bidang vertikal yang membentuk sudut $60^\circ$ membangkitkan rotasi sebesar
$120^\circ$: sebuah $C_3$.
\end{solution}

\begin{solution}{exo:b3:point-groups:11}
Sumbu C=C=C adalah sebuah $C_2$ dan sebuah $S_4$ (putar $90^\circ$, yang
mempertukarkan kedua bidang \ce{CH2}, lalu cerminkan terhadap bidang karbon
pusat); kedua bidang \ce{CH2} adalah $\sigma_d$; dua $C_2$ yang tegak lurus pada
sumbu membagi dua sudut di antara keduanya. Orde 8: $D_{2d}$, akiral. Dalam
\ce{ClHC=C=CHCl}, bidang-bidang, $S_4$, dan $C_2$ aksial dirusak oleh
substitusi itu; hanya satu $C_2$ yang tegak lurus pada sumbu yang tersisa:
$C_2$, kiral --- sebuah alena yang kiral secara aksial.
\end{solution}

\begin{solution}{exo:b3:point-groups:12}
\ce{SF6}: $O_h$ ($h = 48$). \ce{SF5Cl}: $C_4$ melalui Cl dan empat $\sigma_v$:
$C_{4v}$ ($h = 8$). trans-\ce{SF4Cl2}: $D_{4h}$ ($h = 16$). cis-\ce{SF4Cl2}: $C_{2v}$ ($h
= 4$). Masing-masing mempertahankan sebagian operasi $O_h$, yang tertutup
terhadap hasil kali, dan 8, 16, 4 membagi 48. Polar: \ce{SF5Cl} dan
cis-\ce{SF4Cl2}.
\end{solution}

\begin{solution}{pb:b3:point-groups:1}
\textbf{1.} Sepanjang keempat diagonal ruang yang melalui C dan setiap H;
masing-masing memberikan $C_3$ dan $C_3^2$: 8 rotasi.
\textbf{2.} Sepanjang $x$, $y$, $z$ melalui pusat-pusat sisi: $C_2(z)$ memetakan
$(1,1,1)
\to (-1,-1,1)$, sebuah hidrogen. Rotasi sebesar $90^\circ$ terhadap $z$
yang diikuti pencerminan terhadap $xy$ memetakan
$(1,1,1) \to (-1,1,1) \to (-1,1,-1)$, sebuah hidrogen: sebuah $S_4$; bersama
$S_4^3$, 6 operasi.
\textbf{3.} Bidang-bidang $x = \pm y$, $y = \pm z$, $x = \pm z$; $x = y$ memuat
$(1,1,1)$ dan $(-1,-1,1)$: setiap bidang memuat dua ikatan C--H.
\textbf{4.} $1 + 8 + 3 + 6 + 6 = 24$.
\textbf{5.} Inversi memetakan $(1,1,1)$ ke $(-1,-1,-1)$, yang bukan hidrogen:
tidak ada $i$.
\textbf{6.} $E$; $8C_3$; $3C_2$; $6S_4$; $6\sigma_d$.
\textbf{7.} \ce{CH3Cl}: $C_{3v}$, $h = 6$.
\textbf{8.} \ce{CH2Cl2}: $C_{2v}$, $h = 4$.
\textbf{9.} \ce{CHCl3}: $C_{3v}$; \ce{CCl4}: $T_d$, $h = 24$.
\textbf{10.} \ce{CH2FCl}: $C_s$ ($h = 2$); \ce{CHFClBr}: $C_1$ ($h = 1$).
\textbf{11.} Masing-masing mempertahankan operasi-operasi $T_d$ yang
menghormati substitusinya; 6, 4, 6, 24, 2, 1 semuanya membagi 24.
\textbf{12.} Ketiga hidrogen, yang dipermutasikan oleh $C_3$, $C_3^2$ dan ketiga
$\sigma_v$.
\textbf{13.} Semua kecuali \ce{CCl4}: sepanjang sumbu $C_3$ untuk \ce{CH3Cl} dan
\ce{CHCl3}, sumbu $C_2$ untuk \ce{CH2Cl2}, pada \omterm{def:b3:point-groups:mirror}{bidang cermin} untuk
\ce{CH2FCl}, dan tanpa arah yang ditentukan simetri untuk \ce{CHFClBr}.
\textbf{14.} Hanya \ce{CHFClBr} ($C_1$, tanpa $S_n$).
\textbf{15.} Dua enantiomer.
\textbf{16.} Kedua hidrogennya ekuivalen: bidang yang melalui C, F, dan Cl
serta membagi dua H--C--H adalah \omterm{def:b3:point-groups:mirror}{bidang cermin}.
\textbf{17.} Tidak: empat substituen yang berbeda hanya menyisakan $E$; tanpa
$S_n$ molekul itu kiral (dan dapat bersifat polar).
\textbf{18.} $(x, y, z)$: $T_2$.
\textbf{19.} Dua hidrogen adalah ujung-ujung sebuah rusuk tetrahedron; ke-24
operasi membawa rusuk mana pun ke rusuk lain mana pun (6 rusuk itu membentuk
satu himpunan): satu \ce{CH2Cl2}.
\textbf{20.} \ce{CH2FCl}: satu. \ce{CHFClBr}: dua, sepasang enantiomer.
\textbf{21.} Tiga: orto ($C_{2v}$), meta ($C_{2v}$), para ($D_{2h}$).
\textbf{22.} Tiga: 1,2,3- ($C_{2v}$), 1,2,4- ($C_s$), 1,3,5- ($D_{3h}$).
\textbf{23.} Dua (atom-atom Cl bersebelahan atau berseberangan). Hanya satu
\ce{CH2Cl2} yang dikenal, yang menyingkirkan kemungkinan karbon datar dan
mendukung karbon tetrahedral yang diusulkan pada tahun 1874.
\textbf{24.} $1^2 + 1^2 + 2^2 + 3^2 + 3^2 = 24$: \textbf{orde $T_d$, \omterm{def:b3:point-groups:point-group}{grup titik}
metana, adalah $h = 24$}.
\end{solution}
